% part: intuitionistic-logic
% chapter: propositions-as-types
% section: reduction

\documentclass[../../../include/open-logic-section]{subfiles}

\begin{document}

\olfileid[pt]{int}{pty}{red}

\olsection{Redução}

Em !!{derivation}s de dedução natural, !!a{introduction} seguida por
!!a{elimination} constitui um desvio que pode ser ``reduzido''. Por
exemplo, uma $\Intro{\land}$ seguida por uma $\Elim{\land}$ se
``cancelam'', pois a conclusão da $\Elim{\land}$ é sempre um dos
conjuntos da conjunção produzida por \Intro{\land}, e as premissas da
$\Intro{\land}$ são ambos os conjuntos. Assim, se precisamos de
!!a{derivation} de qualquer um dos conjuntos, podemos simplesmente usar
essa !!{derivation}, sem introduzir a conjunção para depois eliminá-la.
Temos
\begin{gather*}
  \bottomAlignProof
  \AxiomC{}
  \RightLabel{$\delta_1$}
  \DeduceC{$!A_1$}
  \AxiomC{}
  \RightLabel{$\delta_2$}
  \DeduceC{$!A_2$}
  \RightLabel{$\Intro{\land}$}
  \BinaryInfC{$!A_1 \land !A_2$}
  \RightLabel{$\Elim{\land}_i$}
  \UnaryInfC{$!A_i$}
  \DisplayProof\bottomAlignProof
  \quad
  \redone
  \quad
  \AxiomC{}
  \RightLabel{$\delta_i$}
  \DeduceC{$!A_i$}
  \DisplayProof
\end{gather*}

Essa simplificação sugere uma relação de redução semelhante nos termos
de prova associados. Se $N_1$ é o termo de prova de~$\delta_1$ e $N_2$
o termo de prova de~$\delta_2$, podemos dizer que o termo de prova da
prova à esquerda se reduz (em um passo) ao termo de prova da prova à
direita:
\[
  \proj{i}{\pair{N_1, N_2}} \redone N_i
\]
No cálculo lambda tipado, essa é a regra de \emph{redução beta} para o
tipo produto.

Na dedução natural, um par de inferências como uma \Intro{\land}
seguida por uma \Elim{\land}, isto é, um par sujeito a redução, é
chamado de \emph{corte}. No cálculo lambda tipado, o termo à esquerda de
$\redone$ é chamado de \emph{redex}, e o termo à direita, de
\emph{reduto}. Ao contrário do cálculo lambda não tipado, no qual apenas
$(\lambd[x][N])Q$ é considerado um redex, no cálculo lambda tipado a
sintaxe é estendida aos termos que envolvem $\pair{N}{M}$,
$\proj{i}{N}$, $\inj{i}{!A}{N}$, $\dcase{N}{x_1}{M_1}{x_2}{M_2}$ e
$\abort{!A}{N}$, com os redexes correspondentes.

De modo semelhante, temos a redução para a disjunção:
\begin{gather*}
  \bottomAlignProof
  \AxiomC{}
  \RightLabel{$\delta$}
  \DeduceC{$!A_i$}
  \RightLabel{$\Intro{\lor}$}
  \UnaryInfC{$!A_1 \lor !A_2$}
  \AxiomC{$\Discharge{!A_1}{x_1}$}
  \RightLabel{$\delta_1$}
  \DeduceC{$!C$}
  \AxiomC{$\Discharge{!A_2}{x_2}$}
  \RightLabel{$\delta_2$}
  \DeduceC{$!C$}
  \DischargeRule{$\Elim{\lor}$}{x_1,x_2}
  \TrinaryInfC{$!C$}
  \DisplayProof\bottomAlignProof
  \quad
  \redone
  \quad
  \AxiomC{}
  \RightLabel{$\delta$}
  \DeduceC{$!A_i$}
  \RightLabel{$\delta_i$}
  \DeduceC{$!C$}
  \DisplayProof
\end{gather*}
Isso corresponde a uma redução nos termos de prova:
\begin{gather*}
  \dcase{\inj[!A_i]{i}{M}}{x_1}{N_1}{x_2}{N_2} \redone \Subst{N_i}{M}{x_i}
\end{gather*}
onde $M$ é o termo de prova de~$\delta$ e $N_i$ é o termo de prova
de~$\delta_i$. Essa é a regra de redução beta para os tipos soma.
Aqui, $\Subst{M}{N_i}{x}$ significa substituir todas as ocorrências
livres da variável~$x$ em $M$ por~$N_i$.

A redução do tipo função corresponde à remoção de um desvio formado por
uma $\Intro\lif$ seguida por uma $\Elim\lif$.
\begin{gather*}
  \bottomAlignProof
  \AxiomC{$\Discharge{!A}{x}$}
  \RightLabel{$\delta$}
  \DeduceC{$!B$}
  \DischargeRule{$\Intro{\lif}$}{x}
  \UnaryInfC{$!A \lif !B$}
  \AxiomC{}
  \RightLabel{$\delta'$}
  \DeduceC{$!A$}
  \RightLabel{$\Elim{\lif}$}
  \BinaryInfC{$!B$}
  \DisplayProof\bottomAlignProof
  \quad
  \redone
  \quad
  \AxiomC{}
  \RightLabel{$\delta'$}
  \DeduceC{$!A$}
  \RightLabel{$\delta$}
  \DeduceC{$!B$}
  \DisplayProof
\end{gather*}
Para termos de prova, isso equivale à redução beta usual:
\begin{gather*}
  (\lambd[\typeof{x}{!A}][N]M)
  \redone \Subst{N}{M}{x}
\end{gather*}

Além das conversões de redução em !!{proof}s, podemos considerar as
conversões de permutação. Elas se aplicam a (sub)!!{proof}s que terminam
em uma \Elim{\lor} seguida por outra !!{elimination}, por exemplo:
\begin{multline*}
  \bottomAlignProof
  \AxiomC{}
  \RightLabel{$\delta$}
  \DeduceC{$!A_i$}
  \RightLabel{$\Intro{\lor}$}
  \UnaryInfC{$!A_1 \lor !A_2$}
  \AxiomC{$\Discharge{!A_1}{x_1}$}
  \RightLabel{$\delta_1$}
  \DeduceC{$!C_1 \land !C_2$}
  \AxiomC{$\Discharge{!A_2}{x_2}$}
  \RightLabel{$\delta_2$}
  \DeduceC{$!C_1 \land !C_2$}
  \DischargeRule{$\Elim{\lor}$}{x_1,x_2}
  \TrinaryInfC{$!C_1 \land !C_2$}
  \RightLabel{\Elim{\land}[i]}
  \UnaryInfC{$!C_i$}
  \DisplayProof\bottomAlignProof
  \quad
  \redone
  \\
  \AxiomC{}
  \RightLabel{$\delta$}
  \DeduceC{$!A_i$}
  \RightLabel{$\Intro{\lor}$}
  \UnaryInfC{$!A_1 \lor !A_2$}
  \AxiomC{$\Discharge{!A_1}{x_1}$}
  \RightLabel{$\delta_1$}
  \DeduceC{$!C_1 \land !C_2$}
  \RightLabel{\Elim{\land}[i]}
  \UnaryInfC{$!C_i$}
  \AxiomC{$\Discharge{!A_2}{x_2}$}
  \RightLabel{$\delta_2$}
  \DeduceC{$!C_1 \land !C_2$}
  \RightLabel{\Elim{\land}[i]}
  \UnaryInfC{$!C_i$}
  \DischargeRule{$\Elim{\lor}$}{x_1,x_2}
  \TrinaryInfC{$!C_i$}
  \DisplayProof
\end{multline*}
Se os termos de prova de $\delta$, $\delta_1$ e $\delta_2$ são $M$,
$N_1$ e $N_2$, respectivamente, a regra de redução correspondente para
os termos de prova, isto é, os termos~$\lambd$, é:
\[
  \proj{i}{\dcase{M}{x_1}{N_1}{x_2}{N_2}} \redone 
\dcase{M}{x_1}{\proj{i}{N_1}}{x_2}{\proj{i}{N_2}}
\]

Naturalmente, podemos eliminar desvios não apenas no final de
!!{proof}s, mas também no interior de !!{proof}s, aplicando a conversão
de redução a uma sub-!!{proof} que termina em um desvio. De modo
semelhante, a redução~$\beta$ aplica-se aos subtermos dos termos~$\lambd$.
Se $N$ é um subtermo de $M$, $N \redone N'$, podemos substituir o
subtermo $N$ em~$M$ por $N'$ para obter $M'$. Nesse caso, também
escrevemos $M \redone M'$. Se existe uma sequência $M = M_1 \redone M_2
\redone M_2 \redone \dots \redone M_k = M'$ (incluindo o caso em que
$k=1$, isto é, $M = M'$), escrevemos $M \red M'$.

Uma !!{proof} à qual nenhuma conversão pode ser aplicada em nenhuma
sub-!!{proof} é chamada de \emph{normal}. De modo semelhante, um
termo~$\lambd$ ao qual nenhuma conversão pode ser aplicada (a nenhum
subtermo) é chamado de \emph{normal}. Um termo~$\lambd$ normal também é
chamado de \emph{forma normal}.

\end{document}
